\section{8/31/2023 Meeting}

Members present: Chris, Braden, Ryan, and Collin

\subsection{Team Goals}
We came up with a list of team goals to pursue this year. These are more general
desires than criteria we will use to measure our success.
\begin{itemize}
      \item Expand the club.
      \item Win worlds.
      \item Greatly improve the notebook.
      \item Improve our engineering design process.
      \item Have a good time and learn a lot.
      \item Go to more competitions and events
      \item Improve funding.
\end{itemize}

\subsection{Iteration \#1 Design Goals}
\begin{itemize}
      \item Be able to intake and launch triballs.
      \item Climb over the central pole easily.
      \item Be able to remove triballs from the goals.
      \item Have both bots be in complement to one another in terms of
            capabilities.
      \item We agreed to, at the moment, ignore the issue of autonomously
            crossing the central pipe as well as climbing to direct our efforts
            to core functionality.
\end{itemize}
We will consider these design goals when we approach the testing phase of our
design, gathering statistics relevant to each point to measure our success.


\subsection{Iteration \#1 Drive Train Brainstorming Analysis}
With these design goals in mind, we considered the following drive mechanisms.
\insertmatrix{design-matrices/drive-iteration-1.tex}{First iteration drive train design matrix.}
\par Note that the speed for the X drive is 1.4. This is because the as the
drive can be viewed as two tank drives joined together at 90\degree{}. At which
point, when both run independently, the velocity vectors of each would add
together, resulting in a forward velocity of $\sqrt{2}$ times faster than tank
with the same gearing and wheel diameter \cite{bib:mecandomniforceanalysis}. The
speed of the other drive trains is mostly dependent on gearing and wheel
diameter, and are assumed equivalent. \par Climbing scores are ranks. Bogie
being the obvious first choice. Next is a tie between tank and mecanum which do
not have any apparent advantages/disadvantages when it comes to climbing. Last
is a tie between X and H which would have to contend with awkward wheel
orientation and a perpendicular wheel respectively. \par Torque scores are
derived similarly to speed: assumed equivalent except for X with analogous
reasoning. \par The scores for simplicity are ranks. Tank is clearly the most
straightforward design, with no special concerns. Next is mecanum, where the
main concern is each wheel being independently-powered (and the additional
width). Then H which requires an additional wheel in the center of the robot.
Then X which requires each wheel be independently-powered and offset by
45\degree{}. Finally, bogie would require two additional center wheels, each
wheel be independently-powered, and an additional linkage. \par Strafing is
simply considered as a boolean value.

\subsection{Iteration \#1 Triball Manipulator Analysis}
Additionally, we considered the following triball manipulator designs.
\insertmatrix{design-matrices/triball-manipulator-iteration-1.tex}{First iteration triball design matrix.}
\par The values for speed (and all categories) are rankings. Flywheel is
fastest, as it can allow continual launching. Catapult and sling rely on a
similar mechanism (tension) and are considered equal in this category. In last
is a claw which would require driving to move the triballs. \par The smallest
design would be a sling, because of its linear movement. Followed by a catapult,
since it moves in an arc. The flywheel design is next which, while it rotates
in-place, would involve the use of large wheels. The largest would be the claw
as it would near certainly involve a deployment mechanism. \par Accessibility
refers to how open the design is to maintenance. All of these designs require at
least occasional servicing (the flywheel may need new motors; catapult, sling,
and claw may need more bands), so this is worthy of consideration. Catapult
comes first in this regard, since the gearing can easily be placed toward the
outside of the robot; plus the fact it can be manipulated to get a better angle.
Flywheel and sling are next as they both lack inherently-inaccessible features.
The claw, comes last, again due to the cumbersome deployment system that would
be a necessity for such a design. \par For simplicity, catapult and sling are
tied for best since they rely on very similar mechanisms (tension) and since
Chris and Collin especially have experience with catapults. Next is the
flywheel, since there could potentially be more moving parts and the design
could involve tuning gearing, velocity control, and compression for consistency.
In last, again, is a claw since it would near-certainly involve a deployment
mechanism.

\subsection{Iteration \#1 Intake Analysis}
For the intake, we considered two different approaches:
\begin{itemize}
      \item A series of compliant wheels.
      \item A series of banded cogs.
\end{itemize}
While we could have put together a design matrix for the intake, we thought it
was best to simply implement both ideas and compare them experimentally.

\subsection{Iteration \#1 Intake Arm Analysis}
Because of the rules involving match loading, we will require the use of an arm
to move the intake. There are two approaches that we considered:
\begin{itemize}
      \item Four bar
      \item Chain bar
\end{itemize}
Again, we decided against a decision matrix for this element of the design;
opting to instead try both approaches to see the superior option.

\subsection{Miscellaneous Ideas/Notes}
\begin{itemize}
      \item Ryan brought up adding a blocker to the 24" robot for opponent
            triballs.
      \item We agreed that we will not consider the pole climbing at this time
            to get a better image of the other subsystems.
      \item We want to use aluminum stock rather the VEX metal due to its
            versatility, cost, and weight. We'll need to find a supplier though,
            Chris asked some teams on the VRC unofficial discord server and
            found a few companies that might do.
      \item Last year Braden began learning the process of anodizing; we hope to
            be able to anodize our robots.
      \item For autonomous functionality, we can simply have two parallel
            odometers and still be able to cross the pipe. Braden brought up
            that we should consider making these metal, as the plastic ones from
            last year (though easy to produce) were fragile.
\end{itemize}

\subsection{Iteration \#1 Design Decisions}
After our analysis, we came to the following design decisions. For both robots,
it should be relatively trivial to add skirts and wings while providing
considerable defensive advantages \cite{bib:blrswikidefense}. We will also
develop odometers for use with both robots, using third-party encoders and a
metal body. We will not consider the endgame with this iteration, at least until
all other aspects of the design have been completed. CAD work has already begun
and we hope to be finished with modeling in the next few weeks.
\subsubsection{24" Iteration \#1 Design Decisions}
For crossing the central barrier, climbing is especially important for both
drives. For this reason, the H and X drives are out of the picture. The bogie
drive would be able to easily cross the central pole, but would be highly
complex. Complexity is a considerable factor for the early season, since club
engagement by new members is heavily influenced by getting robots up and
running; additionally it would provide further practice time for Chris and
testing time for Braden. Since strafing is more of a "bonus" (until competitive
play and testing proves otherwise) we decided to go with the simplest option the
\textit{tank} drive. Later in the season, we will reconsider the use of a
mecanum drive. \par As for the intake, since the 24" has fewer concerns with
space-usage, the \textit{flywheel} design offers great speed with decent
accessibility and simplicity. \par Because of the additional space levied by the
24", we will implement a blocker to prevent opposing teams from scoring via
launching triballs.
\subsubsection{15" Iteration \#1 Design Decisions}
For similar reasons to the 24" bot, the obvious drive train of choice for the
15" is \textit{tank}. Additionally, smart space-usage is critically important
for the 15" and tank drive is the best option in that regard. \par Again, since
size constraints are so strict with the 15" a claw or flywheel would be quite
the burden. While the sling should be more compact than a catapult, Chris and
Collin have experience from Spin Up with fitting a catapult onto a 15" robot.
Since mechanisms relying on tension will inevitably accrue maintenance, the
increased accessibility of the catapult is alluring. For these reasons, we will
presently go forward with a \textit{catapult} design for the 15" robot. We will
later reconsider the use of a sling.

\subsection{Potential Events/Competitions}
From last year, we learned that the best motivator for getting robots completed
was to sign up for a competition. Therefore, we would like to register for a
competition/event in December or early January. \par Prospective ideas include:
\begin{itemize}
      \item Russellville skills at HS competitions in December and January.
      \item Bryant skills at HS competition in January.
      \item CSM VEXU Tournament in Maryland on January 19th.
      \item Houston's competition in March.
\end{itemize}


\insertimage{images/august/8-31-2023-whiteboard-1.jpg}{0.1}
\insertimage{images/august/8-31-2023-whiteboard-2.jpg}{0.1}